\documentclass[11pt]{article}
\usepackage[margin=1in]{geometry}
\usepackage{amsmath,amssymb,amsthm,mathtools}
\usepackage[hidelinks]{hyperref}
\newtheorem{theorem}{Theorem}
\newtheorem{lemma}[theorem]{Lemma}
\newtheorem{corollary}[theorem]{Corollary}
\DeclareMathOperator{\Span}{Span}
\DeclareMathOperator{\Spec}{Spec}
\DeclareMathOperator{\ord}{ord}
\newcommand{\C}{\mathbb C}
\newcommand{\Q}{\mathbb Q}
\newcommand{\Qb}{\overline{\mathbb Q}}
\newcommand{\Lfac}{\mathcal L}
\title{Internal verification of Long's EIT dependency and binary Factorial proof}
\author{Research audit; the results and strategies are credited to their sources}
\date{16 November 2026}
\begin{document}
\maketitle
\noindent\textbf{Finding and scope.}
The main-text argument of Long's binary Factorial manuscript and its pinned
EIT I dependency passes the proof audit below. This is a written internal
verification against published inputs, not a claim of peer review,
proof-assistant verification, or discovery of these results. EIT I is by
Christopher D. Long and Antoine-Auguste Le Blanc; the FC(2) manuscript is
by Christopher D. Long. Exact source versions are in the
\href{source_ledger.json}{source ledger}. The auxiliary finite checks do
not prove analytic continuation or universal nonvanishing.

The FC source is commit
\href{https://github.com/long-mathematics/factorial-conjecture-two-variables/tree/8b1eb46ff530122ed25f588b78bf9fdc98cb9577}{8b1eb46}.
Its EIT input is explicitly pinned to
\href{https://github.com/long-mathematics/exponential-integral-theorem/blob/22d98c4f27e606b5b25fa673c33f9c56c53cc86b/papers/01-exponential-integral-theorem/exponential-integral-theorem.tex}{22d98c4},
not an unspecified later companion theorem. Neither manuscript's
appendices or computational suites are logical inputs to its main result.
The necessary external results are Beukers' linear lifting theorem,
the small-map intermediate-extension theorem, Riemann--Hilbert and
Fourier equivalence for holonomic modules, Pakovich--Muzychuk's rational
permutation-module theorem, and Lindemann--Weierstrass. The applications
and the nonstandard analytic/algebraic steps are checked below.

\section{EIT I: the complete dependency chain}
For $q\in\Qb[x]$ nonconstant and an algebraic degree-zero cycle
$\Delta=\sum d_\alpha[\alpha]$, let
$I(z)=\int_\Gamma e^{zq(x)}dx$, where $\partial\Gamma=\Delta$.
The integrand has an entire primitive in $x$, so only the boundary matters.
We verify the source's assertion
\[
 I(1)\in\Qb\quad\Longrightarrow\quad I\equiv0.               \tag{E}
\]
For $\Delta=[b]-[a]$, $a\ne b$, the value $I(0)=b-a$ then gives
$\int_a^b e^{\xi q(x)}dx\ne0$ for every nonzero algebraic $\xi$, by
replacing $q$ with $\xi q$. This is precisely the FC manuscript's input.

\subsection{Moment coordinates and arithmetic}
Put $n=\deg q\ge2$, $V=\C[x]_{<n-1}$, and divide, for $0\le j<n-1$,
\[
 qx^j=h_jq'+r_j,\qquad \deg r_j<n-1.
\]
Let $C_0:x^j\mapsto r_j$, $D_0:x^j\mapsto-h_j'$, and $C=C_0^T$,
$D=D_0^T$. The degree of $h_j$ is $j+1$, its leading coefficient is
$1/n$, and hence $\Spec(D)=\{-1/n,\ldots,-(n-1)/n\}$.
For $U_j=\int_\Gamma x^j e^{zq}dx$, integration by parts gives
\[
 U'=(C+D/z)U+z^{-1}\sum_{\lambda\in\Lambda} b_\lambda e^{\lambda z},
 \quad \Lambda=\{0\}\cup q(\operatorname{supp}\Delta),
 \quad b_\lambda=\sum_{q(\alpha)=\lambda}d_\alpha h(\alpha).   \tag{6}
\]
All coordinates are $E$-functions: their exponential-series coefficients
are polynomial moments, obtained by evaluating primitives with denominators
at most $nm+j+1$. In a fixed number field, conjugate absolute values grow
at most exponentially in $m$, and one common denominator is
$d^{(n+1)m+j+2}\operatorname{lcm}(1,\ldots,nm+j+1)$, for a common
integrality denominator $d$. Its logarithm is $O(m)$. Equation (6), with
the finitely many exponentials adjoined, supplies differential finiteness.
Thus the actual functions, not just a formal connection, meet Beukers'
hypotheses. The augmented system has no finite pole except zero.

\subsection{Why the homogeneous connection is semisimple}
Over regular values of $q$, $q_*\C$ is a finite permutation local system,
so is semisimple. A finite map of smooth curves is proper and small:
every nondense stratum has dimension zero and every fiber dimension zero,
strictly below the source dimension one. The small-map identity therefore
identifies $Rq_*\C[1]$ with the intermediate extension of that local
system. The relevant published pinpoint is
\href{https://arxiv.org/abs/0712.0349}{de Cataldo--Migliorini, Remark 4.2.4}.
Intermediate extension preserves the direct-sum decomposition into simple
local systems, and Riemann--Hilbert gives a semisimple regular holonomic
$q_+\mathcal O$.

Fourier transform is a Weyl-algebra automorphism and preserves simplicity.
Localization of a simple module at nonzero polynomials in the Fourier
coordinate is either zero or irreducible: the torsion is a differential
submodule, and a nonzero localized submodule meets the original one after
clearing a denominator. The constant direct-image summand transforms to
a module supported at zero and disappears on generic localization.

For the graph generator use $t\mapsto-\partial_\tau$,
$\partial_t\mapsto\tau$. Its relations become
$(\partial_\tau+q)u=0$ and $(\partial_x+\tau q')u=0$.
The coefficient differential in $x$ is therefore $\partial_x-\tau q'$,
not $\partial_x+\tau q'$. Integrating and then putting $\tau=-z$ gives
\[
 M_q=\C(z)[x]/(\partial_x+zq')\C(z)[x],\qquad
 \nabla_z[p]=[\partial_zp+qp].                              \tag{7}
\]
The commutator is zero. Degree comparison proves that $1,\ldots,x^{n-2}$
is a basis of the quotient. Division above gives matrix $C_0+D_0/z$;
the period system has the transpose matrix in (6). This checks the
source's sign and dual conventions and proves semisimplicity for the
homogeneous period connection and its dual.

\subsection{Every rational horizontal projector is constant}
A rational matrix satisfying $P'=[C+D/z,P]$ has no pole away from zero,
by comparing orders. A pole $z^{-m}P_{-m}$ at zero would require
$-mP_{-m}=[D,P_{-m}]$. Eigenvalue differences of $D$ lie in $(-1,1)$,
excluding every negative integer. Thus $P$ is polynomial.

To rule out positive degree, work with the transposed operators $C_0,D_0$.
At each critical point $\tau$ put $m_\tau=\ord_\tau q'$,
$\lambda_\tau=q(\tau)$, and $g_{\tau,k}=q'/(x-\tau)^k$,
$1\le k\le m_\tau$. These form a basis of $V$ by the local Chinese
remainder decomposition. Multiplication by $q$ on $\C[x]/(q')$ acts on
each such vector by $\lambda_\tau$; in particular $C_0$ is semisimple,
even when critical points are multiple or critical values coincide.

Compress $D_0$ to the $\lambda$-eigenspace of $C_0$, denoting the result
by $\bar D_\lambda$. The flag with successive vectors
$g_{\tau,1},\ldots,g_{\tau,m_\tau}$, for $q(\tau)=\lambda$, is preserved.
Indeed
\[
 D_0g_{\tau,k}=-\left(\frac{q-\lambda}{(x-\tau)^k}\right)',
\]
whose leading term at $\tau$ is $-(m_\tau+1-k)a(x-\tau)^{m_\tau-k}$
when $q-\lambda=a(x-\tau)^{m_\tau+1}+\cdots$. At another critical point
over the same value it vanishes modulo that local factor of $q'$.
The diagonal entries of the compressed flag are therefore
$-m_\tau/(m_\tau+1),\ldots,-1/(m_\tau+1)$, all in $(-1,0)$.
If $Q=P^T=Q_dz^d+Q_{d-1}z^{d-1}+\cdots$, $d>0$, the two highest equations are
\[
 [Q_d,C_0]=0,\qquad
 dQ_d=[Q_{d-1},C_0]+[Q_d,D_0].                              \tag{8}
\]
In a nonzero diagonal $\lambda$-block the first commutator in the second
equation vanishes, and the other commutator has spectrum in $(-1,1)$.
It cannot have the positive integer eigenvalue $d$. This proves constancy.
Semisimplicity now supplies a horizontal idempotent for any submodule;
its constancy makes every submodule a constant subspace invariant under
both $C$ and $D$. No generic-critical-value assumption has entered.

\subsection{Algebraic values force exponential polynomials}
Beukers' \href{https://arxiv.org/html/math/0405549}{Theorem 3.2 and Lemma 3.1}
lift every algebraic linear relation at an ordinary nonzero algebraic
point to a polynomial-coefficient functional relation with the prescribed
specialization. Linear independence of all chosen coordinates is not
required. Applying this at one to $U_0(1)-a=0$, with $a$ algebraic, gives
a polynomial row $\phi$ with $\phi(1)=e_0$ and
$\phi U\in E=\Span_{\C(z)}\{e^{\lambda z}:\lambda\in\Lambda\}$.

The set $R=\{\psi:\psi U\in E\}$ is a differential submodule under
$\psi\mapsto\psi'+\psi(C+D/z)$, by (6). The previous subsection gives
$R=W\otimes\C(z)$ for a constant invariant row space $W$.
Every constant annihilator of $W$ annihilates $\phi$ and hence
$\phi(1)=e_0$; thus $e_0\in W$ and $U_0\in E$.

Choose a constant basis $S$ of $W$. In the unique rational-exponential
expansion $SU=\sum r_\lambda(z)e^{\lambda z}$, coefficients satisfy
\[
 r_\lambda'=(C_W-\lambda+D_W/z)r_\lambda+Sb_\lambda/z.        \tag{9}
\]
There are no poles away from zero. At zero a negative integer exponent
would be an eigenvalue of $D_W$, impossible since its spectrum is a subset
of $\Spec(D)\subset(-1,0)$. Thus all coefficients are polynomial.
They are algebraic: a finite number of Taylor coefficients of $U_0$,
all algebraic, determines them by the invertible confluent Vandermonde
matrix for the distinct algebraic $\lambda$. Therefore $I=U_0$ is an
exponential polynomial with algebraic polynomial coefficients.

\subsection{Compact chains supply no nonzero exponential polynomial here}
Choose an embedded finite star through the critical and endpoint values
of $q$. Its inverse graph is connected by transitive polynomial monodromy,
so any degree-zero boundary cycle can be represented on that graph.
Over an open target edge $e$ its pushforward density is an analytic sum
$\psi_e(t)=\sum_i d_{e,i}/q'(x_i(t))$. Endpoint singularities have powers
$1/k-1>-1$ and are integrable.

If a nonzero polynomial $H$ satisfies
$\int_\Gamma q(x)^mH(q(x))dx=0$ for every $m\ge0$, the Cauchy transform
\[
 C_H(s)=\sum_e\int_e\frac{H(t)\psi_e(t)}{s-t}\,dt
\]
vanishes near infinity. The complement of the tree is connected, hence
$C_H$ vanishes there. Its two lateral boundary values at an interior edge
point differ by $\pm2\pi i H(t)\psi_e(t)$. The other edges are holomorphic
locally. Hence every $\psi_e$ vanishes away from the finite zeros of $H$,
and then identically. The original exponential period is consequently zero.

A constant-coefficient differential operator annihilating an exponential
polynomial $I$ gives exactly such a nonzero $H$ by differentiation under
the compact integral. This proves (E) for $n\ge2$. For linear $q=ax+b$,
$I(z)=(az)^{-1}\sum d_\alpha e^{zq(\alpha)}$; Lindemann--Weierstrass and
$\sum d_\alpha=0$ prove (E) directly. This completes the EIT dependency,
including the case of repeated critical points and values. Its extension
to arbitrary polynomial amplitudes is false: for $q=x(x-1)$,
$(2x^2-x+1)e^q=(xe^q)'$ and the integral from zero to one is one.

\section{The nonflat FC(2) proof}
For a nonzero binary polynomial of total degree $D$, write
\[
 P(t,r)=f(rt,r(1-t))=A(t)r^D+A_1(t)r^{D-1}+\cdots,
 \quad d=\deg A\le D.
\]
Restriction of homogeneous forms to $x+y=1$ is injective, so $A\ne0$.
We first assume $d>0$. An all-moment counterexample can be specialized
to algebraic coefficients without changing $D,d$: impose the rational
moment equations, the linear degree constraints, and an inverse for the
leading coefficient of $A$. Their ideal is proper because of the given
complex point; Noetherianity and Nullstellensatz give an algebraic point.
No preservation of monodromy under specialization is needed.

\subsection{Exact kernel and the correction that cannot be discarded}
Set $M_n=\Lfac(f^n)/\Gamma(Dn+2)$ and
\[
 H(w)=\frac1w\sum_{n\ge0}M_nw^{-n}
 =\int_0^1 C(w,t)dt,
 \quad C(w,t)=\frac1w\int_0^\infty e^{-r}r
       E_{D,2}(P(t,r)/w)dr .                               \tag{10}
\]
The bound $|E_{D,2}(z)|\le e^{|z|^{1/D}}$ justifies the integral and
summation for sufficiently large $|w|$, locally uniformly on bounded
$t$-sets. All positive moments zero means $H=1/w$.

On $V=\{A\ne0\}$ let
$\Phi_t(r)=(P/A)^{1/D}=r+p_1+O(r^{-1})$, $p_1=A_1/(DA)$,
and let $\chi(t,Y)=Y-p_1+O(Y^{-1})$ be its exterior inverse. Define
\[
 F(t,Y)=e^{Y-\chi}\chi\chi_Y/Y=\sum_{m\ge0}F_m(t)Y^{-m},
 \qquad F_0=e^{p_1}=:E(t).
\]
The root filter
$E_{D,2}(X^D)=(DX)^{-1}\sum_{j=0}^{D-1}\omega^{-j}e^{\omega^jX}$
has the additional subtraction $X^{-1}$ when $D=1$.
Split the radial integral into a compact part and a uniformly exterior
part before applying this filter. For $\rho^D=A(t)/w$, the distinguished
singular summand on the exterior part is
\[
 \frac{\rho^{D-1}}{DA(t)}\int e^{-(1-\rho)Y}F(t,Y)dY .       \tag{11}
\]
Here $1/(\rho w)=\rho^{D-1}/A$, which checks the prefactor.
On compact subsets of $V$, $|F_m|\le C R_0^m$ by the convergent exterior
Laurent expansion. For $R_1>R_0$, the elementary integrals
$J_m(\eta)=\int_{R_1}^\infty e^{-\eta Y}Y^{-m}dY$ satisfy
\[
 J_0=e^{-\eta R_1}/\eta,\quad J_1=-\log\eta+\text{entire},\quad
 J_{m+1}=(e^{-\eta R_1}R_1^{-m}-\eta J_m)/m.
\]
Their logarithmic coefficients are $(-1)^m\eta^{m-1}/(m-1)!$.
After removing them, the entire remainders obey
$R_1^m\|H_{m+1}\|\le e^{LR_1}/m+(LR_1/m)R_1^{m-1}\|H_m\|$
on $|\eta|\le L$. This gives a uniform bound and normal convergence
after multiplication by $F_m$. Thus (11) has a simple pole at
$w=A(t)$ with residue $E(t)$ and a logarithmic coefficient
\[
 B(w,t,\rho)=\frac{\rho^{D-1}}{DA(t)}
   \sum_{m\ge1}\frac{(-1)^mF_m(t)}{(m-1)!}(1-\rho)^{m-1}.    \tag{12}
\]
It is holomorphic on the finite cover $\rho^D=A(t)/w$, $w\ne0$, $t\in V$.
The only finite exceptional $w$-values for the full kernel are $0,A(t)$.
At $A(t)=0$, the original integral is normally convergent on compact
$w$-sets away from zero: its integrand is bounded by
$C r\exp(-r/2+C'(1+r)^{1-1/D})$ after shrinking the $t$-neighborhood.
These estimates check the parameter uniformity and continuation used below.

Choose a target star through $0,A(0),A(1)$ and all critical values,
with a regular center. Its inverse graph is a tree, and its path from
the actual endpoint zero to the actual endpoint one replaces $[0,1]$
in (10). Each arm $v$ has a signed sheet vector $n_v$.
At a nonzero leaf the kernel gives a corrected germ
\[
 L_v(w)=\sum_i n_{v,i}\left(
     E(t_i(w))t_i'(w)+\int_{c_i}^{t_i(w)}B(w,t,\rho)dt\right),\qquad
 c_i=\lim_{w\to v}t_i(w).                                 \tag{13}
\]
For rational $H$, this whole germ vanishes. To check the correction sign,
write the local contribution as
$\int_0^\epsilon[a(x)/(z-x)+b(z,x)\log(z-x)]dx$.
Integration by parts with $q(z,x)=\int_0^x b(z,s)ds$ gives
$\int[a(x)+q(z,x)]/(z-x)dx$ plus a holomorphic boundary term.
For each Puiseux class $-1<\delta\le0$, division of its analytic coefficient
by $x-z$ leaves $h_\delta(z,z)\int x^\delta/(z-x)dx$.
For $\delta<0$ its nonholomorphic term is a nonzero multiple of
$h_\delta(z,z)z^\delta$; for $\delta=0$ it is $h_0(z,z)\log z$.
Finite-power monodromy first kills fractional powers and detects the log;
then the distinct monodromy eigenvalues separate all fractional classes.
Meromorphy therefore forces every coefficient to vanish, proving (13).
Discarding the integral correction and asserting a relation for the pole
density alone at this stage would be invalid.

\subsection{Transport, with histories and root sheets retained}
Continue (13) along paths in the base minus zero and critical values.
Each integral retains its fixed lower endpoint and appends the path of
its moving upper endpoint. It is lifted to (12)'s cover with $\rho=1$
at the upper endpoint; lower-endpoint roots are never reset after a loop.
On a small base disk choose $s(w)^D=w$. The cover is the product of that
disk with $\{(t,\sigma):A(t)=\sigma^D\ne0\}$, via $\rho=\sigma/s(w)$.
The differential in (12) is holomorphic on this product. Fixed compact
history pieces therefore integrate holomorphically in $w$.

Approach a nonzero critical value $\lambda$. A short terminal piece has
a joint holomorphic primitive $J(w,t)$, and its moving endpoint is
$c+O((w-\lambda)^{1/k})$. It supplies only nonnegative Puiseux exponents.
Thus every negative exponent of the continued zero identity comes solely
from $E(t_i(w))t_i'(w)$. For any monodromy permutation $g$, continuation
along a realizing loop gives annihilation of $g n_v$ by every such
negative-coefficient functional. Histories may depend on more than $g$;
only their negative coefficients factor through the permutation action.

\subsection{The top rational component and simultaneous unweighted traces}
Use the geometric full-cycle labeling of the sheets. For every invariant
periodic subspace $V_k$, $k\mid d$, put
$U_d=\bigcap_{k<d}V_k^\perp$. The rational irreducibility of this space
is the exact top-component case of
\href{https://arxiv.org/html/0710.4085v2}{Pakovich--Muzychuk, Theorem 3.1
and Section 4.1}. It is rational, not absolute complex irreducibility.
Its application does not assume that $A$ is indecomposable.

The signed path module $M=\sum_v\Q[G]n_v$ contains $U_d$.
Here are the checks on the geometric input. At each center the path has
one incoming and one outgoing edge, so $\sum_vn_v=0$ and the nonzero leaves
already generate $M$. Local-monodromy averaging of $n_v$ gives the uniform
sheet vector at each endpoint, or their difference if endpoint values
coincide. The embedded stars formed by an endpoint's lifted arm and the
disjoint rays to infinity separate their boundary labels. Their convex
hulls are disjoint when the endpoint values agree and meet at most at
one boundary label otherwise. Consequently the two averages differ in
the first case and cannot both have root-of-unity mean zero in the second.
The vector $(1,\zeta,\ldots,\zeta^{d-1})$ belongs to $U_d\otimes\C$ and
detects $M$. Rational orthogonal projectors onto $M,U_d$ commute: both
centralize the full cycle, whose centralizer is the commutative circulant
algebra. Irreducibility then gives $U_d\subset M$. This is also the
path-module statement in Pakovich--Muzychuk, Proposition 4.1, with the
actual endpoint and zero-leaf conventions checked here.

The unweighted inverse-derivative vector
$u=(1/A'(t_i(w)))_i$ belongs to $U_d\otimes\C$.
Indeed, a proper invariant block comes from a polynomial decomposition
$A=P\circ R$, $\deg R\ge2$ (the intermediate field has a unique place
over infinity, so L\"uroth generators can be chosen polynomial).
On a block $R(c)=\xi$, $P(\xi)=w$, its coordinate sum is
$P'(\xi)^{-1}\sum_{R(c)=\xi}1/R'(c)=0$, by the coefficient of $t^{-1}$
in $1/(R(t)-\xi)=O(t^{-\deg R})$ at infinity. These are constant rational
linear equations, valid coefficientwise in every Puiseux expansion.
All proper decompositions are enforced simultaneously.

\subsection{Arithmetic contradiction and the endpoint case}
If $\lambda\ne0$ is critical, let $K$ be the largest local degree over it
and $\beta\ne0$ the coefficient vector of $u$ at exponent $1/K-1$.
It is algebraic and belongs to $U_d\otimes\C$.
No smaller local degree or positive-order term of the analytic weight
contributes to that exponent. The corresponding weighted vector is
$\phi_i=\beta_i\exp(A_1(c_i)/(D\lambda))$. Transport makes it orthogonal
to $U_d\otimes\C$, which is defined over $\Q$ and contains
$\overline\beta$. Hence
\[
 0=\langle\phi,\overline\beta\rangle
   =\sum_{\beta_i\ne0}|\beta_i|^2e^{A_1(c_i)/(D\lambda)}.    \tag{14}
\]
Grouping equal algebraic exponents gives strictly positive algebraic
coefficients. Lindemann--Weierstrass contradicts (14). This conjugate
pairing does not require real coefficients of $A$.

Thus all critical values are zero. Counting multiplicities of roots of
$A$ and $A'$ shows $A=a(t-c)^d$. Every nonzero endpoint value is regular,
so at its leaf (13) evaluates to the signed value $E(p)/A'(p)$; the
shrinking integral is zero. Different endpoint values give one nonzero
uncancelled value. Equal endpoint values give a nonzero common $v$ and
would require $e^{(A_1(0)-A_1(1))/(Dv)}=A'(0)/A'(1)$.
Hermite--Lindemann forces the exponent difference to vanish, then forces
$A'(0)=A'(1)$. But these derivatives are $dv/(-c)$ and $dv/(1-c)$,
which differ. For $d=1$ endpoint values cannot agree. This contradiction
completes the nonflat proof in all degrees, not just a sampled range.

\section{Flat case, full binary conclusion and exact Gaussian consequences}
If $A$ is constant, normalize it to one and specialize algebraically.
The multinomial radial expansion gives, for $K=\sum k_j$, $J=\sum jk_j$,
\[
 C_{n,k}=\frac{n!}{(n-K)!}\frac{\Gamma(Dn-J+2)}{\Gamma(Dn+2)},
 \qquad 0\le C_{n,k}\le1,
\]
with zero assigned when $K>n$. For fixed $k$ it is asymptotic to
$D^{-J}n^{K-J}$. The summable bound $\prod_j\|A_j\|^{k_j}/k_j!$ justifies
uniform dominated convergence and proves
\[
 \lim_{n\to\infty}\frac{\Lfac(f^n)}{\Gamma(Dn+2)}
     =\int_0^1 e^{A_1(t)/D}dt\ne0.                          \tag{15}
\]
EIT supplies the last inequality for arbitrary algebraic $A_1$; a constant
$A_1$ is immediate. Together with the nonflat proof and specialization,
this verifies FC(2). If a tail of moments were zero, a sufficiently high
power of $f$ would have every positive moment zero. Therefore
\[
 \{f\in\C[y_1,y_2]:\Lfac(f^m)=0\text{ eventually}\}=\{0\}.  \tag{16}
\]
This is an internal reproduction of Long's conclusion, not a new priority claim.
The proof does not assert a nonzero limit in (15) for every transcendental
coefficient choice: $f=(x+y)^2+4\pi i x$ has $A_1(t)/2=2\pi i t$, so that
limit is zero although $\Lfac(f)=6+4\pi i\ne0$. Algebraic specialization
is applied to a putative all-moment counterexample, rather than to its
individual analytic limit.

\begin{corollary}[Exact two-block radial classification]
Let $P=f((X_1^2+X_2^2)/2,(X_3^2+X_4^2)/2)$ over a characteristic-zero
field. If its pure Gaussian moments vanish eventually, then $f=0$.
Consequently all-marker eventual safety in this entire invariant
subalgebra is equivalent to $P=0$.
\end{corollary}
\begin{proof}
The exact factorial/Gaussian identity transfers the hypothesis to (16).
The conclusion extends from $\C$ to any characteristic-zero field by
embedding the finitely generated coefficient field into $\C$.
The substitution is injective, for example by restricting to
$X_2=X_4=0$ and comparing monomials in $X_1^2,X_3^2$.
All-marker safety includes the unit marker; zero is trivially safe.
\end{proof}

\begin{corollary}[Every commutative coefficient $\Q$-algebra]
Let $B$ be any commutative $\Q$-algebra, including a nonreduced or
non-Noetherian one. For $f\in B[y_1,y_2]$, the following are equivalent:
\begin{enumerate}
\item $\Lfac_B(f^m)=0$ for all sufficiently large $m$;
\item every coefficient of $f$ is nilpotent in $B$;
\item $f$ is nilpotent in $B[y_1,y_2]$;
\item $\Lfac_B(qf^m)=0$ eventually for every fixed polynomial $q$.
\end{enumerate}
The same classification holds in the two-block Gaussian invariant
subalgebra, retaining every original Gaussian polynomial marker.
\end{corollary}
\begin{proof}
Reduce (1) modulo any prime ideal of $B$ and pass to the fraction field.
The field result makes every coefficient zero there, hence every coefficient
belongs to the nilradical of $B$. A polynomial has only finitely many
coefficients; their finitely generated nil ideal is nilpotent, by the
pigeonhole bound $1+\sum_i(e_i-1)$ if their nilpotence exponents are $e_i$.
Thus (2) implies (3) without a Noetherian hypothesis. Nilpotence gives (4),
and the unit marker gives (1). The Gaussian pure identity gives the same
necessity, and nilpotence gives vanishing against all original markers.
\end{proof}
In particular, reduced coefficient algebras have only the zero safe radial
phase. Eventual vanishing cannot be replaced by first-moment vanishing;
over nonreduced rings the coefficient ideal need not be zero.
For example, if $\epsilon^2=0$ but $\epsilon\ne0$, the polynomial
$\epsilon(y_1-y_2)$ is nonzero and every positive pure moment is zero.
Thus even the all-positive-moment hypothesis does not imply literal zero
over a nonreduced coefficient algebra.

There is also an exact ordinary normalization ideal, without a reduced-ring
assumption. Let $J\subset B$ be generated by the coefficients of $f$, and
let $N_f$ consist of scalars $u$ for which every actual marked period
$u\sum_m\Lfac_B(qf^m)t^m/m!$ is polynomial in $t$. Then
\[
 N_f=\Gamma_J(B)=\bigcup_{k\ge1}\operatorname{Ann}_B(J^k).    \tag{17}
\]
To prove necessity, reduce modulo $\operatorname{Ann}_B(u)$: multiplication
by $u$ embeds this quotient into $B$, so eventual vanishing of the scaled
unit period gives eventual vanishing in that quotient. The preceding
corollary makes its finitely many coefficients nilpotent, hence $J^k u=0$
for some $k$. Conversely $J^k u=0$ gives $uf^m=0$ for all $m\ge k$,
uniformly over all markers. The proof also covers $u=0$.
For example, in $B=\Q[s,t]/(s^2t^2)$ and $f=sy_1$, the ideal is exactly
$(t^2)$, rather than $(t)$. This follows immediately by the monomial basis
of the quotient. Thus this classification retains the ordinary torsion
that a reduced safe-locus computation would lose.

\section{Boundary of the audit}
General FC in three or more variables, general SFC, and absence of
pure-null polynomial functions on triangles remain unproved here.
The binary result does not prove weighted Gamma FC for arbitrary shapes,
or a universal Gaussian Mathieu statement. The unrestricted Image
counterexample remains valid. The current EIT repository also has Paper VI,
\emph{Exponential Periods and Factorial Moments}, commit
\href{https://github.com/long-mathematics/exponential-integral-theorem/tree/92db79eedeb064694ff20e9201d73fc2bc0857f4}{92db79e};
its higher-dimensional family assertions require a separate audit and do
not supply unrestricted FC. Nothing here upgrades the manuscripts to
peer-reviewed status or asserts formal verification of the analytic proofs.
\end{document}
